\chapter*{Over het Open Logic Project}
\addcontentsline{toc}{chapter}{Over het Open Logic Project}


De \textit{Open Logic Text} is een opensourceleerboek dat gezamenlijk wordt
geschreven. Het behandelt formele metalogica en formele methoden en begint op
een niveau waarvoor een inleidende cursus formele logica als voorkennis wordt
verondersteld. Hoewel het boek bestemd is voor een niet-wiskundig publiek, in
het bijzonder voor studenten filosofie en informatica, is de behandeling
wiskundig nauwkeurig.

Enkele onderwerpen die al in de tekst staan, worden mogelijk nog niet volledig
behandeld, en veel paragrafen moeten nog grondig worden herzien. We zijn van plan
de tekst in de toekomst met meer onderwerpen uit te breiden. Ook zijn we van
plan de tekst uit te breiden met onder meer een woordenlijst, een
literatuurlijst, historische toelichtingen, afbeeldingen en betere uitleg,
evenals paragrafen die de relevantie van de resultaten voor filosofie,
informatica en wiskunde toelichten, en meer opgaven en voorbeelden. Vindt u een
fout of hebt u een suggestie,
\href{https://github.com/OpenLogicProject/OpenLogic/wiki/Contributing}{laat het projectteam dat dan weten}.

Het project werkt volgens de uitgangspunten van open source. De tekst is niet
alleen vrij beschikbaar; we stellen ook de LaTeX-broncode beschikbaar onder de
Creative Commons Naamsvermelding-licentie. Deze licentie geeft iedereen het
recht ons werk te downloaden, te gebruiken, te wijzigen, anders te ordenen,
om te zetten en opnieuw te verspreiden, mits de makers naar behoren worden
vermeld. Meer informatie vindt u op de website van het Open Logic Project:
\href{http://openlogicproject.org/}{openlogicproject.org}.
